\begin{explanationbox}
	\textbf{\textcolor{CExplanation}{一句话直觉}}
	球面距离公式来源于球心角的余弦定理，将经纬度转化为向量夹角。

	\medskip
	\textbf{\textcolor{CExplanation}{核心拆解}}
	\begin{description}[style=nextline, leftmargin=2.8em, labelsep=0.5em, itemsep=0.45em]
		\item[\textbf{要点一（球心角计算）：}]
			先由经纬度求出球心角 $\theta$ 的余弦：
			\[
				\cos\theta = \sin\varphi_1\sin\varphi_2 + \cos\varphi_1\cos\varphi_2\cos(\Delta\lambda).
			\]
		\item[\textbf{要点二（弧长转化）：}]
			球面距离等于半径乘以球心角的弧度：$l = R\theta$。
		\item[\textbf{要点三（坐标对应）：}]
			若将点转化为空间直角坐标，则 $\cos\theta$ 是两点单位向量的点积。
	\end{description}

	\medskip
	\textbf{\textcolor{CExplanation}{几何本质（球心角余弦）}}
	球心角 $\theta$ 对应大圆上弧 $AB$ 所对圆心角，其余弦由纬度、经度差通过球面余弦定理给出。

	\medskip
	\textbf{\textcolor{CExplanation}{代数意义（向量点积）}}
	设 $A,B$ 对应的空间直角坐标为单位向量，则 $\cos\theta = \overrightarrow{OA}\cdot\overrightarrow{OB}$，公式是该点积在球坐标下的展开。

	\medskip
	\textbf{\textcolor{CExplanation}{考点价值（高考应用）}}
	\begin{itemize}[leftmargin=*, itemsep=0.5em, topsep=0.4em]
		\item \textbf{考法一（直接套用）：} 给出经纬度求球面距离，直接代入公式。
		\item \textbf{考法二（逆向求角）：} 已知距离反求球心角或经纬度。
	\end{itemize}

	\medskip
	\textbf{\textcolor{CExplanation}{顿悟点}}
	记住 $\cos\theta$ 的表达式：正弦乘积 + 余弦乘积 × 经度差余弦。结构类似两角差的余弦公式，但多了纬度正弦项。

	\medskip
	\textbf{\textcolor{CExplanation}{使用场景}}
	\begin{itemize}[leftmargin=*, itemsep=0.5em, topsep=0.4em]
		\item \textbf{场景一（地球问题）：} 已知地球表面两点经纬度，求最短距离。
		\item \textbf{场景二（立体几何）：} 球面上两点距离转化为球心角求解。
	\end{itemize}

\end{explanationbox}
